\documentclass[11pt]{article}
\usepackage[margin=1in]{geometry}
\usepackage{amsmath,amssymb,amsthm,hyperref}
\newtheorem{theorem}{Theorem}\newtheorem{lemma}[theorem]{Lemma}
\newtheorem{proposition}[theorem]{Proposition}
\newcommand{\R}{\mathbb R}\newcommand{\F}{\mathbb F}\newcommand{\E}{\mathbb E}
\newcommand{\one}{\mathbf1}\newcommand{\relu}{\operatorname{ReLU}}
\title{Superpolynomial strong-prior ReLU mistakes\protect\\with margins only on observed inputs}
\author{ChatGPT-6 Astra\thanks{Independent AI research. Sole author. The rail atom theorem, path-cost construction, and conditional Fourier argument used below are credited to the preceding jointly authored paper.}}
\date{21 November 2026; working proof draft}
\hypersetup{hidelinks,pdftitle={Superpolynomial strong-prior ReLU mistakes with margins only on observed inputs},pdfauthor={ChatGPT-6 Astra}}
\emergencystretch=3em
\begin{document}\maketitle
\begin{abstract}
We construct a Boolean target of polynomial one-hidden-layer ReLU
complexity and a fixed order of distinct Boolean inputs on which an
isotropic Gaussian ReLU learner makes superpolynomially many confident
mistakes in the strong-prior limit. Training conditions only on positive
margins at observed inputs; there is no initial restriction on unobserved
outputs. For a parameter $k$, the physical input dimension is $O(k)$,
the target has a Boolean circuit of size $O(k)$ and a shallow ReLU teacher
of width $O(k^2)$ with polynomially bounded coefficients, and
$2^{k-1}$ designated predictions are wrong with posterior probability
tending to one. The result holds at every sufficiently large fixed
learner width. The calibration prefix and sufficient learner width are
large. We also obtain a separate fixed-positive-temperature width-first
ReLU consequence and a standard-softplus strong-prior extension. The proof adds a control rail to the existing unequal-rail parity
construction. An explicit dual certificate has the correct sign on
observed examples and vanishes on all unobserved examples. A conditional
Fourier argument excludes every off-code competitor and identifies a
unique wrong minimum-energy completion. This replaces the earlier
full-cube absolute-margin hypothesis. The proof is conventional and has
not received independent mathematical review.
\end{abstract}

\section{The learner and the quantitative conclusion}
On $X=\{-1,1\}^d$, let $\bar x=(1,x)$ and
\[
 F_\theta(x)=\sum_{j=1}^N a_j\relu(U_j\cdot\bar x),\qquad
 \relu(t)=\max(t,0).
\]
Independently, the coordinates of $U_j$ have law
$N(0,\alpha\sigma^2)$ and the readouts have law
$N(0,\beta\sigma^2)$, where $\alpha,\beta>0$. At each fixed $d,N$
these constants may depend on $d,N$ but not on $\sigma$.
After a history $H$ of labeled inputs, condition the Gaussian prior on
\[
                   f(x)F_\theta(x)>1/2\qquad(x\in H).       \tag{1}
\]
There are no other output constraints. The strong-prior limit is
$\sigma\downarrow0$ at fixed $d,N$.

\begin{theorem}[Observed-margin-only mistakes]\label{thm:main}
For every integer $k\ge2$, put
\[
 m=6k+4,\quad \ell=8k+4,\quad R=64,\quad
 d=\ell(2R+4)+2R+2=1056k+658,
\]
\[
 n_{\max}=544k+338,\quad N_0=2^{m+n_{\max}}=2^{550k+342},
 \qquad C_k=3(m+2)(d+1),\quad T_k=2^{k-1}.
\]
There is a Boolean function $f_k:X\to\{-1,1\}$ and a fixed permutation
of $X$ with $T_k$ designated rounds such that, for every finite
$N\ge N_0$, the correct-label posterior probability at every designated
round tends to zero under (1).

The target has a bounded-fan-in Boolean circuit of size $O(d)$ and a
one-hidden-layer ReLU teacher of width at most $C_k$, output margin one,
and coefficients bounded in absolute value by $C_k$. Thus the number
of designated errors is exponential in the input dimension and in the
displayed circuit-size bound, and is
$\exp(\Omega(\sqrt{C_k}))$ in the displayed teacher width.

For every $\eta>0$, at all sufficiently small positive $\sigma$,
posterior-label majority and the sign of the posterior mean score are
wrong at all designated rounds, and fresh
posterior sampling makes a mistake at all of them with probability at
least $1-\eta$. The same input order works for every eligible learner
width and every pair of positive layer-variance constants.
\end{theorem}

The sequence is deliberately unrepresentative. The theorem is not an
iid lower bound. The calibration prefix is exponential in dimension,
and $N_0$ is a very large sufficient width. The result is not asserted
at every width just above the teacher width.

\section{The inherited rail and path ingredients}
We use results from the preceding jointly authored paper \cite{shared}.
Their hypotheses and the parts needed here are stated explicitly.
The new argument starts with the observed-label control dual in
Section~\ref{sec:control}; the atom and path constructions are not new here.

For positive even $n$ define
\[
 \kappa_n=\frac{\binom{n-2}{(n-2)/2}}{2^{n-1}},\qquad
 a(n)=\frac{\sqrt n}{\kappa_n}.
\]
Partition the physical coordinates into disjoint rail pairs
$(B_j^0,B_j^1)$, each rail having even size in $[R,R+\Delta]$, with
$R\ge6\Delta$. A full transversal selects one whole rail from every
pair. If $A_t$ is such a transversal, the exact rail atom inequality is
\[
 \sum_t a(|A_t|)|\widehat{\relu(U\cdot\bar x)}(A_t)|
                                                    \le\|U\|_2. \tag{2}
\]
Equality for a nonzero atom holds precisely on the balanced,
zero-bias rays supported on one full transversal: all nonzero incoming
coordinates have equal absolute value. Such an atom has Fourier support
only on subsets of its transversal. Therefore the ReLU variation norm
\[
 V(F)=\inf\left\{\sum_j|c_j|\|u_j\|_2:
                    F=\sum_jc_j\relu(u_j\cdot\bar x)\right\}
\]
satisfies
\[
 V\left(\sum_t c_t\chi_{A_t}\right)
                          =\sum_t a(|A_t|)|c_t|.          \tag{3}
\]
Signed coordinate-flip averaging on its active coordinates attains
the cost of each parity with at most $2^{|A_t|}$ balanced rays.
The elementary bounds
\[
 n\le a(n)\le3n,\qquad a(n+2)-a(n)\ge2                  \tag{4}
\]
follow from the central-binomial formula; for the second use
$a(n+2)/a(n)=\sqrt{n(n+2)}/(n-1)\ge(n+2)/n$.

In $G=\F_2^m$, take $P=4k+3$ prefix coordinates and a path on
$2k+1$ further coordinates. There is one singleton column $a_j$ for
every coordinate and one column for each adjacent path pair, giving
$\ell=8k+4$ old columns. Every singleton has increment four.
Successive edges have increments $6,2,6,2,\ldots,6,2$.
Every short rail has size $R$ and its long rail has size $R+\delta_j$.

Add one new rail pair of sizes $R$ and $R+2$, with independent syndrome
sign $\tau$. For the old pairs put
$(-1)^{\xi_j}=\chi_{B_j^0\cup B_j^1}(x)$ and
$z=\sum_j a_j\xi_j\in G$. Let $b(x)$ be the parity of all short rails,
including the new one. The retained full-transversal characters are
\[
 b(x)\chi_s(z),\qquad b(x)\tau\chi_s(z),\qquad s\in G,
 \quad \chi_s(z)=(-1)^{s\cdot z}.
\]
Their sizes and weights are
\[
 n_s=n_0+\sum_j\delta_j(a_j\cdot s),\quad n_0=(\ell+1)R,
 \qquad w_s=a(n_s),\quad w'_s=a(n_s+2),\quad w_c=w'_0.     \tag{5}
\]
All rails satisfy (2) with $\Delta=6$. The variables $z,\tau,b$ are
independent uniform syndromes: the old singleton columns span $G$,
and the two nonempty rail parities within each pair are independent.

Write $M=2^m$ and
$\widehat w(u)=M^{-1}\sum_s\chi_s(u)w_s$.
The path construction supplies $U=v_0+H\subseteq G$ with
\[
 |U|=K=2^{k-1},\quad 0\notin U,\quad
 (U\cup\{0\})\cap((U\cup\{0\})+a_j)=\varnothing
                          \quad\hbox{for every old column }a_j. \tag{6}
\]
Concretely, fix all prefix bits to one, all even-position path bits to
one, the terminal path bit to zero, and the parity of the other $k$
odd-position path bits to zero. The $k-1$ free bits give $H$.
No old column has more than two nonzero entries, so none belongs to
$U$; the path independence then gives (6). The costs satisfy
\[
 \widehat w(u)<0\ (u\ne0),\qquad
 \widehat w\hbox{ is strictly ordered on }U,\qquad
 W:=\max_{u\in U}|\widehat w(u)|
 \le\frac{8\sqrt{2\pi}}{n_0-1}
                         \left(\frac2{R-1}\right)^P.     \tag{7}
\]
The extra short rail only multiplies the old Laplace integrand by
$e^{-Rt}$, preserving positivity, strict order, and this bound.
The inherited proof uses
$a(n)=\sqrt{2\pi}(n-3/4)-\int_0^\infty e^{-nt}\rho(t)dt$,
where $0<\rho(t)\le8\sqrt{2\pi}e^t$. The path's triangular transfer
matrices give a signed-digit expansion with ratio in $(0,1/3)$ at
every $t>0$. Integration preserves the strict order; it is not inferred
from numerical cancellation. The order can be specified without real
comparisons. If $\eta_1,\ldots,\eta_k$ are the free odd-position bits
with their prescribed total parity, put
$\epsilon_i=(-1)^{\eta_i+1}$ and
$d_j=\prod_{i=k-j+1}^k\epsilon_i$ for $1\le j<k$.
Increasing cost is lexicographic order of $(d_1,\ldots,d_{k-1})$
with $+1$ before $-1$, because the cost is the negative integral of
the path's positive-factor signed-digit expansion.

We also use the conditional Fourier observation from \cite{shared}.
Let $\mathcal D$ be the downward span of retained transversals and
$\mathcal C=b\{g(z,\tau)\}$ their complete code span.
For $A\in\mathcal D\cap\mathcal C^\perp$, assign each physical character
to its first defective rail pair, where its support is neither of the
two complete rails. Write $A=\sum_j A_j$, with old pairs before the new
pair. Conditional on $(z,\tau)$ the pieces are orthogonal, and
\[
 \ell_j(z,\tau):=\E[A_j^2\mid z,\tau]
 \quad\hbox{is invariant under translation by that pair's column}. \tag{8}
\]
For an old pair this sends $(z,\tau)$ to $(z+a_j,\tau)$; for the new
pair it flips $\tau$. Indeed two physical characters have nonzero
conditional covariance only if their symmetric difference is a union
of full rail pairs representing a code frequency. At a defective pair,
complementation inside the whole pair cannot remain a subset of one
rail. Their restrictions must agree there. This proves first-defect
orthogonality and the translation invariance.

\section{A target and calibration using ordinary labels}
Define
\[
 t(z,\tau)=
 \begin{cases}+1&\tau=+1\text{ or }z=0,\\
              -1&\tau=-1\text{ and }z\ne0,\end{cases}
 \qquad f(x)=b(x)t(z(x),\tau(x)).                         \tag{9}
\]
First reveal every input with $\tau=+1$, and every input with
$\tau=-1$ and $z\notin U$. Then reveal whole minus fibers of $U$ in
increasing order of $\widehat w(u)$. At a designated stage let
$L\subseteq U$ be the nonempty remaining suffix and $v$ its unique
cheapest element. Consider
\[
 F_v(x)=b(x)[g_v(z(x))+\tau(x)/2],\qquad
                         g_v=\one_{\{0\}}+\one_{\{v\}}.   \tag{10}
\]
It has correct margin at least $1/2$ on every observed input.
On the unobserved fiber $(v,-1)$ it is $b/2$, while the target is $-b$.
Its exact variation cost is
\[
                   C_v=\widehat w(0)+\widehat w(v)+w_c/2, \tag{11}
\]
because $\widehat g_v(s)=(1+\chi_s(v))/M\ge0$ and (3) applies.

\section{An observed-label dual certificate}\label{sec:control}
Let
\[
 E=H^\perp\setminus(H+\operatorname{span}(v_0))^\perp.
\]
Then $E$ is nonempty, $0\notin E$, $|E|=M/(2K)$, and every character
in $E$ equals $-1$ on $U$. Put
\[
 c_E=-M\widehat w(v)>0,\quad p_E=c_E/|E|,\qquad
                         \widehat q(s)=w_s-p_E\one_E(s).
\]
Fourier inversion gives
\[
 q(z)=M\widehat w(z)-c_E\one_H(z)+c_E\one_U(z).           \tag{12}
\]
Thus $q(0)>0$, $q<0$ outside $U\cup\{0\}$, and
\[
                  q(u)=M[\widehat w(u)-\widehat w(v)]
                                     \quad(u\in U).     \tag{13}
\]
In particular $q(v)=0$, $q>0$ on $L\setminus\{v\}$, and $q<0$
on the removed prefix. With $\delta=(K+1)/M$, define
\[
 \gamma=\frac{w_c+w_0-2\E_z[q(z)\one_L(z)]}{1-\delta},
\]
\[
 h(z)=-q(z)+2q(z)\one_L(z)
                  +\gamma[1-\one_{\{0\}}(z)-\one_U(z)]. \tag{14}
\]
Expectations here are uniform on $G$.

\begin{lemma}[Admissible dual coefficients]\label{lem:coefficients}
For every $k\ge2$, $\gamma>0$, $0<p_E<2\min_s w_s$, and
\[
 \widehat h(0)=w_c,\qquad |\widehat h(s)|<w'_s\quad(s\ne0). \tag{15}
\]
Also $|\widehat q(s)|\le w_s$, with equality precisely when $s\notin E$.
\end{lemma}
\begin{proof}
By (7), $0\le\E(q\one_L)\le KW$ and $p_E\le2KW$.
Since $\sqrt{2\pi}<3$ and $2/63<1/16$, the loose bound
\[
                     2KW<24\,2^k16^{-(4k+3)}<1/8        \tag{16}
\]
holds for $k\ge2$. Thus $\gamma>0$ and
$\gamma\le(w_c+w_0)/(1-\delta)\le6(n_0+2)/(1-\delta)$.
Since $\delta\le3\,2^{-5k-6}<1/2$,
\[
             \gamma\delta\le36(512k+322)2^{-5k-6}<3/4.  \tag{17}
\]
At $k=2$ the last expression is $48456/65536<3/4$; it decreases
thereafter because its numerator grows by less than two while its
denominator grows by 32. Equations (4) and (16) prove the assertions
for $q$. Averaging (14) gives $\widehat h(0)=w_c$.
For $s\ne0$, use $|\widehat{\one_U}(s)|\le K/M$ to obtain
\[
 |\widehat h(s)|\le|\widehat q(s)|+2KW+\gamma\delta
                                      <w_s+1<w'_s.
\]
Finally $q(0)=M[\widehat w(0)+\widehat w(v)]>0$ since
$\widehat w(0)\ge n_0$ and $W<1$.
\end{proof}

Define the physical dual
\[
                    \Lambda(x)=b(x)[q(z(x))+\tau(x)h(z(x))]. \tag{18}
\]
Its coefficients are $\widehat q(s)$ on main transversals,
$\widehat h(s)$ on mixed transversals, and zero on all other full
transversals. Lemma~\ref{lem:coefficients} and (2) give
\[
                            |\E_x\Lambda F|\le V(F).     \tag{19}
\]
This is a dual for the unrestricted physical ReLU network.

\begin{lemma}[Support, signs, and complementarity]\label{lem:signs}
The dual vanishes on every unobserved input and has the target sign on
every observed input where it is nonzero. Also $F_v=f/2$ wherever
$\Lambda\ne0$, and $\E\Lambda F_v=V(F_v)=C_v$.
\end{lemma}
\begin{proof}
The logical values are
\[
\begin{array}{c|cc}
 z&q+h&q-h\\\hline
 z\notin U\cup\{0\}&\gamma>0&2q-\gamma<0\\
 0&0&2q(0)>0\\
 u\in L&2q(u)\ge0&0\\
 u\in U\setminus L&0&2q(u)<0 .
\end{array}                                                   \tag{20}
\]
These signs agree with (9). Only the minus faces of $L$ are unobserved;
at $v$ both entries vanish. Formula (10) gives the exact margins.
Its main coefficients are nonnegative and vanish on $E$, saturating
$q$; its control coefficient $1/2$ saturates $\widehat h(0)=w_c$.
\end{proof}

Thus $C_v$ is the global minimum cost for the closed observed margins.
We must still exclude another physical minimizer with a different sign
at the query. This step uses no unobserved margin.

\section{Rigidity against every physical ReLU competitor}
\begin{proposition}[The unique minimum-cost output]\label{prop:unique}
Among all finite ReLU representations satisfying the closed observed
margins, the only output with representation cost at most $C_v$ is $F_v$.
\end{proposition}
\begin{proof}
For such a representation of cost $B$,
\[
 B\ge\E\Lambda F\ge\tfrac12\E|\Lambda|
                              =\E\Lambda F_v=C_v.
\]
All inequalities are equalities. Hence $F=f/2$ wherever $\Lambda\ne0$.
Every nonzero neuron saturates the rail atom and coefficient bounds.
The strict mixed inequalities exclude all mixed transversals except
control $(0,1)$; the strict main coefficients exclude $E$; zero dual
coefficients exclude every omitted full transversal. The equality
characterization in (2) puts $F$ in the downward span
$\mathcal D_{\rm ret}$ of
\[
 \{\text{main transversals }s\notin E\}
                         \ \cup\ \{\text{control transversal}\}. \tag{21}
\]
The same is true of $F_v$.

Put $D=F-F_v$. Its possible nonzero logical fibers lie in
\[
 S=\{(0,+1)\}\cup(L\times\{-1\})
       \cup((U\setminus L)\times\{+1\})\cup\{(v,+1)\}.    \tag{22}
\]
Project onto the full code:
$P_{\mathcal C}D=b\E[bD\mid z,\tau]$.
Both the projection and the residual $A=D-P_{\mathcal C}D$ vanish
outside $S$. The residual has the first-defect decomposition in (8).
Conditional orthogonality implies that each $\ell_j$ vanishes outside $S$.

For any old pair, $S$ contains at most one endpoint of its column
translation, by (6). Translation invariance forces $\ell_j=0$
everywhere. For the tag pair, $\ell_{\rm tag}$ is invariant under
flipping $\tau$. The only $z$ with both tag faces in $S$ is $v$.
Therefore $A_{\rm tag}$ is supported on $z=v$.

A character assigned to the tag defect is complete on every old pair.
Under (21), its old index is $s\notin E$; the control's old index is
$0\notin E$ too. Fix the physical coordinates of the tag pair.
Then $A_{\rm tag}$, as a function of the old coordinates, is $b_{\rm old}$
times a function of $z$ whose Walsh support excludes $E$.
This function is supported only at $v$. A nonzero function supported
at one point has every Walsh coefficient nonzero. Since $E$ is
nonempty, it must vanish. Thus $A_{\rm tag}=0$ and $F\in\mathcal C$.

No omitted complete transversal is contained in a retained one.
Consequently (21) restricts the code projection to
\[
                   F=b[g(z)+a\tau],\qquad
                           \widehat g(s)=0\quad(s\in E). \tag{23}
\]
At any $z\notin U\cup\{0\}$ both tag faces have active weights:
$g(z)+a=1/2$ and $g(z)-a=-1/2$. Thus $a=1/2$, and $g=0$ at every
such $z$. The active minus face at zero gives $g(0)=1$.
For a remaining $u\ne v$ the active plus face gives $g(u)=0$;
for a removed $u$ the active minus face does the same.
Hence $g=\one_{\{0\}}+t\one_{\{v\}}$.
Any $s\in E$ gives $0=(1-t)/M$, so $t=1$ and $F=F_v$.
\end{proof}

The control face supplies the missing nonnegativity in the reduced
completion problem. Strict mixed coefficients exclude interaction
features in a saturating physical network. The last single-fiber
argument is necessary because column independence fails at the
current query's two tag faces.

\section{Widths, Gaussian conditioning, and teacher size}
Each $F_v$ has exactly $M/2$ nonzero main parity coefficients and one
control coefficient. Every full transversal has at most
$n_{\max}=n_0+4\ell+2=544k+338$ coordinates.
By (3), a minimum-cost representation uses at most
$(M/2+1)2^{n_{\max}}\le M2^{n_{\max}}=N_0$ neurons.
Zero-padding covers every $N\ge N_0$. Balancing gives
\[
 \frac12\sum_j\left(\frac{a_j^2}{\beta}
                          +\frac{\|U_j\|^2}{\alpha}\right)
                   \ge\frac{V(F)}{\sqrt{\alpha\beta}},    \tag{24}
\]
with equality for the balanced minimum-cost representation.
More directly, the left side dominates the cost of its own
representation divided by $\sqrt{\alpha\beta}$, so Proposition
\ref{prop:unique} identifies the output of every energy minimizer.

The exponential width is not solely an artifact of this representation
bound. Every exact minimum-cost network has at least $M/2+1$ nonzero
neurons. Indeed each nonzero neuron must be a balanced ray on a single
retained transversal. Its Fourier support contains no other complete
transversal, whereas $F_v$ has $M/2+1$ distinct nonzero full-transversal
coefficients. Each requires at least one such neuron. Thus this
particular minimum-norm mechanism genuinely uses an exponentially
wider representation than the polynomial-width teacher.

\begin{proposition}[Calibration is also intrinsically large for these completions]
If an observed-margin history has $F_v$ as its unique minimum-variation-
cost output, then it contains at least $M/2+1$ distinct observed inputs.
\end{proposition}
\begin{proof}
More generally, whenever a finite minimum-cost representation exists
for a history of $h$ inputs, some minimum-cost representation has at
most $h$ nonzero neurons. Normalize each nonzero neuron's incoming
vector to length one and absorb its readout sign into its atom.
Its positive coefficient is then its contribution to variation cost.
If there are more than $h$ such atoms, their restrictions to the
observed inputs are linearly dependent: there is a nonzero vector $z$
whose weighted atom sum vanishes on the whole history.
Both sufficiently small coefficient perturbations $c\pm\epsilon z$
remain nonnegative and preserve every observed output. Minimum cost
forces $\sum_jz_j=0$, since otherwise one sign would reduce the cost.
Increase one perturbation until its first coefficient reaches zero.
Cost and all observed outputs are unchanged. Repeating gives at most
$h$ nonzero neurons. If the minimum-cost output is unique, it must still
be $F_v$. The preceding active-width lower bound then yields
$h\ge M/2+1$.
\end{proof}
Thus neither exponential resource is merely a loose upper bound of
this proof: its particular unique completions require exponential
active width and exponentially many earlier observations. This does
not rule out a different superpolynomial-mistake mechanism with a
shorter calibration or smaller learner width.

Designate the first input of each successive minus fiber in increasing
cost order, completing that fiber before proceeding. Calibration and
these fibers give a permutation of the cube. At each designated query
the target is $-b$ and the unique minimizing output is $b/2$.
The closed correct-sign event $f(x)F(x)\ge0$ therefore excludes every
energy minimizer. Coercivity and continuity give a strictly larger
energy infimum on that event, unless it is empty.
Multiplying the readouts of $F_v$ by $1+\epsilon$ makes all observed
margins strict, at arbitrarily small extra energy. An open parameter
ball about this witness lower-bounds the posterior denominator.
Here are the details of the probability comparison. Write $I$ for the
quadratic energy in (24), let $E_0$ be its constrained minimum, and let
$E_1>E_0$ be the correct-sign infimum. Choose $a\in(E_0,E_1)$ and a
strictly feasible open ball $\mathcal B$ with $I<a$.
Choose $\varepsilon\in(0,1)$ with $(1-\varepsilon)E_1>a$.
In parameter dimension $D_N=N(d+2)$, the unnormalized numerator is at most
\[
 e^{-(1-\varepsilon)E_1/\sigma^2}
        \int e^{-\varepsilon I/\sigma^2}\,d\theta
       =c_\varepsilon\sigma^{D_N}
                          e^{-(1-\varepsilon)E_1/\sigma^2},
\]
whereas the denominator is at least
$\operatorname{vol}(\mathcal B)e^{-a/\sigma^2}$.
Their ratio tends to zero. This integrates all parameterizations and
makes no optimizer tie-breaking assumption.

The same argument applies outside any fixed neighborhood of the unique
minimizing output table. Moreover $|F_\theta(x)|\le C_{d,\alpha,\beta}I(\theta)$.
Inserting this polynomial energy factor into the Gaussian tail integral
still gives exponential decay beyond a sufficiently large fixed energy
level. Compact-output concentration and this tail bound imply
$\E[F_\theta(x)\mid H]\to F_v(x)$ as well. Hence the posterior mean
score has the same wrong limiting sign. There are finitely many
designated rounds at fixed $k,N$; a union bound gives the simultaneous
sampling and majority assertions.

Finally a shallow teacher score for (9) is
\[
 F_{\rm teacher}(x)
       =b(x)\left[\sum_{i=1}^m\chi_{e_i}(z(x))+m\tau(x)+1\right]. \tag{25}
\]
For $\tau=+1$ its bracket is at least one. For $\tau=-1$ it is one
at $z=0$ and at most minus one elsewhere.
It is a sum of $m+2$ physical parities. Parity on at most $d$ bits is
represented by the three-hinge hats on its attainable coordinate sums,
using at most $3(d+1)$ neurons, incoming weights zero or one, biases
bounded by $d+1$, and readouts bounded by two.
The logical coefficients in (25) give width $C_k$ and coefficient
bound $C_k$. Computing rail parities, bounded-degree old syndromes,
the control parity, and the final Boolean test uses $O(d)$ Boolean
gates. This completes Theorem~\ref{thm:main}.

\begin{proposition}[A polynomial mistake bound from a finite teacher prior]
There is a finite prior on one-hidden-layer integer-weight ReLU networks
of width $C_k$, independent of the target's particular labels, whose
observed-margin posterior-label majority makes at most
\[
                 C_k(d+2)\log_2(2C_k+1)=O(k^3\log(k+1))
                                                               \tag{25a}
\]
mistakes on every input sequence labeled by $f_k$.
It may be padded to the same learner width as in Theorem~\ref{thm:main}.
No polynomial computation-time bound is claimed for this alternative.
\end{proposition}
\begin{proof}
The three-hinge parity hats above use integer incoming weights, biases,
and readouts. The logical teacher coefficients are integers too.
Thus the unit-margin teacher belongs to the class of width-$C_k$
networks with all $C_k(d+2)$ parameters in
$\{-C_k,\ldots,C_k\}$, padding with zero readouts when necessary.
Give every such parameter table equal prior mass.
There are at most $(2C_k+1)^{C_k(d+2)}$ tables.
The teacher survives every observed margin event (1).
Whenever posterior-label majority makes a mistake, at least half
the surviving mass predicts the wrong sign and is removed by the next
margin observation. Hence after $L$ mistakes the surviving prior mass
is at most $2^{-L}$ and at least the original teacher mass.
Taking logarithms proves (25a).
\end{proof}
This elementary counting argument isolates the role of the prior:
the labels have a small integer-weight teacher, while the Gaussian
minimum-energy rule prefers the much wider low-norm completions.
It does not identify an efficient way to discover the teacher.

\section{Quantitative stability and the size of the energy gap}
The preceding rigidity argument has an explicit norm version.
All norms below use the uniform physical or logical cube, as appropriate.
Let $O=\{(z,\tau):q(z)+\tau h(z)\ne0\}$ be the active observed fibers.

\begin{lemma}[Quantitative observability of the retained space]
For every $R_0\in\mathcal D_{\rm ret}$,
\[
                  \|R_0\|_2^2\le18K\,\E[\one_O R_0^2]. \tag{Q1}
\]
\end{lemma}
\begin{proof}
First, a logical function $g$ with Fourier support excluding $E$ obeys
\[
             \|g\|_2^2\le\frac M{|E|}
                              \E[\one_{\{z\ne v\}}g(z)^2]
                       =2K\,\E[\one_{\{z\ne v\}}g(z)^2]. \tag{Q2}
\]
Indeed Cauchy--Schwarz in its Fourier expansion gives
$|g(v)|^2\le(M-|E|)\|g\|_2^2$; remove the single point's
$M^{-1}|g(v)|^2$ contribution.

Decompose $R_0=G+A$, where $G=b[g(z)+a\tau]$ is its full code
projection, and split $A$ by first defective pairs. These components
are orthogonal even after restriction to the logical set $O$.
For an old defect, column invariance and (6) give
$\|A_j\|_2^2\le2\E[\one_O A_j^2]$.
For the tag defect, apply (Q2) after fixing the physical tag variables:
its old Fourier support excludes $E$. Tag invariance and the fact that
every $z\ne v$ has at least one face in $O$ then give
\[
 \|A_{\rm tag}\|_2^2
       \le2K\,\E[\one_{\{z\ne v\}}A_{\rm tag}^2]
       \le4K\,\E[\one_O A_{\rm tag}^2].
\]

For the code component, both faces are in $O$ outside $U\cup\{0\}$.
Thus $(1-\delta)a^2\le\E[\one_O G^2]$, and $a^2\le2\E[\one_O G^2]$.
Choose one observed active face $\tau_z$ at each $z\ne v$.
Using $g(z)^2\le2[g(z)+a\tau_z]^2+2a^2$ and the normalization
$1/(2M)$ on the two-face cube gives
\[
 \E[\one_{\{z\ne v\}}g^2]\le4\E[\one_O G^2]+2a^2
                                      \le8\E[\one_O G^2].
\]
Equation (Q2) yields
$\|G\|_2^2=\|g\|_2^2+a^2\le18K\,\E[\one_O G^2]$.
Sum the conditionally orthogonal components to prove (Q1).
\end{proof}

For the fixed stage let $\vartheta<1$ be the largest absolute
dual-coefficient/weight ratio among excluded full transversals.
Let $K_{\rm ret}$ be a retained-ray stability constant from
\cite{shared}; one may take
\[
                    K_{\rm ret}=2^{16}(\ell+1)^2R^4.    \tag{Q3}
\]
The full ray stability estimate restricts to any subset by charging
a closest omitted ray's coefficient to the omitted Fourier mass.
All full-transversal weights are at most $6(\ell+1)R$, so the same
elementary restriction bound in \cite{shared} gives (Q3).

For a finite ReLU representation of cost $B$, its dual deficit
$D=B-\langle\Lambda,F\rangle$ controls the off-space error:
\[
       \|F-P_{\mathcal D_{\rm ret}}F\|_2^2
               \le\frac{K_{\rm ret}}{1-\vartheta}\,BD.   \tag{Q4}
\]
To check the coefficient factor, write $Q_{\rm ret}$ for the retained
absolute coefficient sum. Every atom satisfies
\[
 |\langle\Lambda,h\rangle|
 \le Q_{\rm ret}(h)+\vartheta[Q_{\rm full}(h)-Q_{\rm ret}(h)]
 \le\vartheta\|U\|+(1-\vartheta)Q_{\rm ret}(h).
\]
Combine retained-ray stability with weighted Cauchy--Schwarz.
This also proves (Q4) for a norm-attaining representation.

\begin{proposition}[An explicit correct-sign cost gap]\label{prop:quantgap}
Set
\[
 B_{\rm cap}=C_v+1,\quad
 H_{\rm cap}=\sqrt{d+1}\,B_{\rm cap}+3/2,\quad
 \mu=\min_O|q(z)+\tau h(z)|>0,
\]
\[
 A_0=\frac{K_{\rm ret}B_{\rm cap}}{1-\vartheta},\qquad
 A_1=\frac{36K H_{\rm cap}}\mu+(36K+1)A_0.
\]
Every score satisfying the closed observed margins and having
representation cost $B\le B_{\rm cap}$ obeys
\[
                        \|F-F_v\|_2^2\le A_1(B-C_v).     \tag{Q5}
\]
Consequently every correct-sign completion has variation cost at least
\[
                 C_v+\min\{1,\ 2^{-d}/(4A_1)\}.          \tag{Q6}
\]
\end{proposition}
\begin{proof}
Put $S=\langle\Lambda,F\rangle-C_v\ge0$, so $D+S=B-C_v$.
On $O$, $F_v=f/2$ and the observed margin implies
$|F-F_v|=f(F-F_v)$. Thus
$\E[\one_O|F-F_v|]\le S/\mu$ and
$\E[\one_O(F-F_v)^2]\le H_{\rm cap}S/\mu$.
Write $R_0=P_{\mathcal D_{\rm ret}}(F-F_v)$ and
$E_0=F-P_{\mathcal D_{\rm ret}}F$. Equation (Q4) bounds
$\|E_0\|_2^2$ by $A_0D$.
On $O$, $(R_0)^2\le2(F-F_v)^2+2E_0^2$.
Apply (Q1) and then use the global orthogonality of $R_0,E_0$:
\[
 \|F-F_v\|_2^2
 \le36K H_{\rm cap}S/\mu+(36K+1)A_0D
 \le A_1(D+S).
\]
At a correct-sign physical query, $|F-F_v|\ge1/2$ at that point,
so the uniform-cube squared norm is at least $2^{-d}/4$.
This proves (Q6) below the cap, while costs above the cap are separated
by one. Norm attainment on the finite cube extends it to $V(F)$.
\end{proof}

The constants do not conceal a doubly exponential cost separation.
There is an absolute $A>0$ such that, uniformly over the stages,
\[
 b_k:=\min\left\{\min_{u\in U}|\widehat w(u)|,\
                 \min_{\substack{u,v\in U\\u\ne v}}
                     |\widehat w(u)-\widehat w(v)|\right\}
                       \ge e^{-A k\log(k+1)}.           \tag{Q7}
\]
Here is a direct justification from the inherited positive integrals.
Their density $\rho(t)$ extends continuously to zero with
$\rho(0)=7\sqrt{2\pi}/32>0$. This follows from the formulas in
\cite{shared}: $k''(0)=1/48$, $B'(0)=1/8$, and $B''(0)=1/128$.
For two distinct path words, the pointwise gap is at least
$e^{-c k}t^k$ on a fixed small interval. The prefix factor is at least
$(ct)^P$, and all remaining positive path factors are at least
$e^{-ck}$ there. Integrate over $[1/n_0,2/n_0]$ for sufficiently
large $k$ to obtain
$e^{-ck}n_0^{-(P+k+1)}$. For an individual absolute coefficient,
the positive path term selecting no edges is at least
$(ct)^{2k}$, giving
$e^{-ck}n_0^{-(P+2k+1)}$ instead.
Both imply (Q7), since $n_0=O(k)$ and $P=O(k)$.
The finitely many smaller $k$ are absorbed by increasing $A$.

In particular $p_E\ge2Kb_k$, while every nonretained mixed coefficient
has weight gap greater than one. Therefore
\[
 1-\vartheta\ge\frac{2Kb_k}{\max_t a(|A_t|)},\qquad
 \mu\ge\min\{1,2Mb_k\}.
\]
The first inequality uses $p_E<1$ and also covers omitted transversals.
All other quantities in (Q3)--(Q6) grow at most exponentially in $k$
or polynomially in $k$ and $b_k^{-1}$. Thus, for another absolute $A$,
\[
             \inf_{\rm correct}V-C_v\ge e^{-A k\log(k+1)}. \tag{Q8}
\]
This is a quantitative lower bound on the cost gap; the construction
still has a large calibration prefix and exponential sufficient width.

\begin{proposition}[A deterministic minimum-norm consequence]
At each designated stage, every observed-margin-interpolating ReLU
score with variation cost less than the correct-sign threshold in
(Q6) predicts the wrong label. In particular every global minimizer
of the two-layer squared Euclidean parameter cost does so, at every
width $N\ge N_0$, with any positive relative weighting of the two
layers. The same assertion holds for approximate minimizers whose
objective excess is below the corresponding rescaled gap.
\end{proposition}
\begin{proof}
The first assertion is the contrapositive of (Q6). The balancing
identity (24) identifies the globally minimum parameter energy as
$C_v/\sqrt{\alpha\beta}$. A parameter energy less than
$(C_v+g)/\sqrt{\alpha\beta}$ gives representation cost less than
$C_v+g$, where $g$ is the gap in (Q6), and therefore the wrong sign.
\end{proof}
This is a statement about global or sufficiently accurate constrained
norm minimization. It does not assert that gradient descent reaches
these minimizers or that an arbitrary regularized training loss has
the same solutions.

\begin{theorem}[A sufficient finite prior scale]\label{thm:finite-scale}
For the canonical ReLU layer variances
$\alpha=N/d$, $\beta=1/N$, there is an absolute constant $A>0$
such that, for every $k\ge2$, $N\ge N_0$, and $0<\eta<1/2$, the choice
\[
 0<\sigma\le
 \frac{\exp[-A k\log(k+1)]}
      {\sqrt{\,N(1+\log N)+\log(1/\eta)\,}}              \tag{Q9}
\]
makes the sum of the correct-label posterior probabilities under (1)
over all designated rounds at most $\eta$. Thus posterior-label majority is
wrong at every designated round, and fresh posterior sampling is wrong
at all of them with probability at least $1-\eta$.
The absolute constant is not numerically optimized.
\end{theorem}
\begin{proof}
Whiten the canonical parameters as
$U_j=\sqrt{N/d}\,u_j$, $a_j=v_j/\sqrt N$.
Then the energy is $I=\frac12\sum_j(\|u_j\|^2+v_j^2)$ and the
ReLU output is $d^{-1/2}\sum_jv_j\relu(u_j\cdot\bar x)$, independently
of width apart from the number of summands.
At each stage its minimum is $E_0=\sqrt d\,C_v=O(k^{3/2})$.
By (Q8) the correct-sign energy gap is at least
$G=e^{-A_0k\log(k+1)}$, after increasing an absolute $A_0$.

A minimum-cost star has whitened parameter norm polynomial in $k$.
Increase its readouts by $1+t$, where
$t=e^{-A_2k\log(k+1)}$ is small enough that the energy increase is
less than $G/8$. The observed margins then exceed $1/2$ by at least
$t/2$. On a bounded energy ball, the uniform-cube output map is
Lipschitz in the whitened parameters with a constant polynomial in $k$,
independent of the number of summands; this is the Cauchy--Schwarz bound
on the two layer differences. There is consequently a strictly feasible
parameter ball of radius $r\ge e^{-A_3k\log(k+1)}$ with energy at most
$E_0+G/4$.

In dimension $D_N=N(d+2)$ this ball has volume at least
$(2r/\sqrt{D_N})^{D_N}$, using its inscribed cube.
Let $E_1$ be the correct-sign infimum, so $E_1\ge E_0+G$.
Choose $\varepsilon=G/[4(E_0+G)]$ and bound the Gaussian numerator
by its tail at $E_0+G$ rather than at the possibly larger $E_1$.
The probability ratio is at most
\[
 \exp\left\{A_4D_N[k\log(k+1)+\log D_N]
                                  -\frac{G}{2\sigma^2}\right\}
                                                               \tag{Q10}
\]
for $\sigma\le1$. Here the denominator uses the ball, and
$\varepsilon^{-1}$ and $r^{-1}$ are at most
$e^{A_5k\log(k+1)}$.
Since $D_N=O(Nk)$, $N\ge N_0$, and $T_k\le2^k$, increasing the
absolute constant $A$ in (Q9) makes (Q10) at most $\eta/T_k$
uniformly over all stages. Sum these bounds.
\end{proof}

\section{A finite hinge penalty and a Gaussian-prior Gibbs law}
There is nevertheless one ordinary training penalty for which the
same conclusion follows exactly. For the actual observed set $H$, let
\[
 \mathcal L_H(F)=\frac1{|H|}\sum_{x\in H}
                        (1/2-f(x)F(x))_+,
 \qquad n_{\max}=n_0+4\ell+2.
\]
The loss is averaged over the distinct observed examples, including
the calibration prefix; it is not reweighted by the dual certificate.

\begin{theorem}[Finite hinge regularization]\label{thm:hinge}
Fix $k\ge2$, $N\ge N_0$, and positive layer constants
$\alpha,\beta$. Write $I$ for the Gaussian energy in (24).
Choose the finite penalty coefficient
\[
                  \lambda>\frac{12M n_{\max}}{\sqrt{\alpha\beta}}.
                                                               \tag{H1}
\]
At every designated round, every global minimizer of
$I(\theta)+\lambda\mathcal L_H(F_\theta)$ has output $F_v$ and
predicts the wrong label. The Gaussian-prior Gibbs law with density
proportional to
\[
       \exp\{-[I(\theta)+\lambda\mathcal L_H(F_\theta)]/\sigma^2\}
                                                               \tag{H2}
\]
therefore has correct-label probability tending to zero as
$\sigma\downarrow0$ at every designated round.
The coefficient in (H1) may be exponential in $k$. Law (H2) is a
hinge-loss Gibbs law; it is not claimed to equal a normalized logistic
or other specified binary observation model.
\end{theorem}
\begin{proof}
Let $a_*=(\alpha\beta)^{-1/2}$ and define nonnegative empirical
coefficients
\[
 c_x=a_*\frac{|H|}{2^d}|\Lambda(x)|\qquad(x\in H).
\]
The dual sign table and $\frac12\E|\Lambda|=C_v$ give
\begin{align*}
 I+\lambda\mathcal L_H-a_*C_v
 &\ge \frac1{|H|}\sum_{x\in H}
 \bigl[(\lambda-c_x)(1/2-fF)_+
                         +c_x(fF-1/2)_+\bigr].
\end{align*}
Every coefficient $\widehat q(s)$ is positive and at most $w_s$,
so $\|q\|_\infty\le3M n_{\max}$.
By (14), $\|h\|_\infty\le\|q\|_\infty+\gamma$;
the proof of Lemma~\ref{lem:coefficients} gives
$\gamma\le12n_{\max}$.
Since $M\ge2$, this implies
$\|\Lambda\|_\infty\le12M n_{\max}$.
Condition (H1) makes $\lambda-c_x>0$ for every observed example.
The balanced star attains objective $a_*C_v$ with zero hinge loss.
Any minimizer must therefore have zero hinge loss and minimum
parameter energy, so Proposition~\ref{prop:unique} forces $F_v$.

The objective is continuous and coercive, since it dominates $I$.
Its closed correct-sign set has a strictly higher minimum.
A small open parameter ball around a minimizing star has energy-plus-loss
below that minimum. The same finite-dimensional Gaussian integral
comparison as above applies, using the objective's lower bound by $I$
for the tail integral. This proves the Gibbs assertion.
\end{proof}
In particular the finite penalty is a sufficient condition for a
specific global training objective, not a claim about convergence of
an optimization algorithm. For the canonical layer variances its
required order is at most $e^{O(k)}$.

\section{A separate fixed-temperature, width-first ReLU consequence}
The strict variation-cost gaps also have a different asymptotic
consequence. This statement uses mean-field normalization and fixes a
small positive prior scale before taking the width limit.

\begin{theorem}[A fixed positive scale for each constructed target]
For each fixed $k\ge2$, there is $\tau_k>0$ such that the same target,
input permutation, and $T_k$ designated rounds have the following
property. In the model
\[
 F_N=N^{-1}\sum_{j=1}^N a_j\relu(U_j\cdot\bar x),
 \qquad a_j,(U_j)_i\ \text{independent }N(0,\tau^2),
\]
fix any $0<\tau<\tau_k$ and use only the observed margins (1).
The correct-label posterior probability at every designated round
then tends to zero exponentially as $N\to\infty$.
Moreover there is an absolute $A>0$ for which every
$0<\tau\le e^{-A k\log(k+1)}$ satisfies the conclusion.
The sufficient scale can still be extremely small.
\end{theorem}
\begin{proof}
On a finite cube the normalized signed ReLU atoms form a compact
spanning set. They span all output tables because a ReLU can isolate
each cube vertex. Their symmetric convex hull is the compact unit ball
of $V$, and Caratheodory's theorem gives attainment of $V$ by finitely
many atoms. Thus Proposition~\ref{prop:unique} applies to every
variation-norm minimizer, not only to an already chosen representation.
At each stage the closed correct-sign constraint has minimum variation
cost $L_v>C_v$: its intersection with a bounded norm sublevel is compact
and excludes the unique minimizer.

Let $I_\tau$ be the Cramer rate of the one-neuron output vector.
Its cumulant is finite near zero. Gaussian readout integration and
homogeneity show that the closure of its finite domain is
$\{\xi:V^*(\xi)\le\tau^{-2}\}$: inside, use
$|\xi\cdot\Phi|\le V^*(\xi)|a|\|U\|$; outside, an open incoming cone
with larger correlation makes the Gaussian integral diverge.
Readout symmetry makes the cumulant nonnegative. Hence
\[
 I_\tau(z)\le V(z)/\tau^2.
\]
Conversely choose a supporting dual for $V(z)$ and shrink it by
$1-\varepsilon$. In $d+2$ Gaussian coordinates,
$2|a|\|U\|\le a^2+\|U\|^2$ bounds the cumulant by
$\frac{d+2}{2}\log(1/\varepsilon)$. Therefore
\[
 I_\tau(z)\ge (1-\varepsilon)V(z)/\tau^2
                         -\frac{d+2}{2}\log(1/\varepsilon). \tag{26}
\]
Choose $\varepsilon>0$ with $(1-\varepsilon)L_v>C_v$.
For sufficiently small $\tau$, the lower rate bound on the closed
correct-sign event exceeds $C_v/\tau^2$.
The strict history is open and contains arbitrarily slightly amplified
$F_v$, whose rate is at most its variation cost divided by $\tau^2$.
Cramer's upper and lower bounds give a strict conditional exponential
rate gap. There are finitely many stages, so one $\tau_k>0$ works for
all of them. Smaller positive $\tau$ preserve the strict inequalities.
For the stated quantitative choice, use the uniform gap (Q8),
$C_v=O(k)$, and choose
$\varepsilon$ of order $e^{-A_0k\log(k+1)}/(C_v+1)$.
The penalty in (26) is $O(k^2\log(k+1))$.
A sufficiently large absolute exponent in the displayed $\tau$ bound
makes the positive rate gap dominate that penalty at every stage.
\end{proof}

\section{Standard softplus with observed margins}
A further extension uses the canonical width-dependent Gaussian layers.
It is separate from the arbitrary fixed layer scales allowed for ReLU.

\begin{theorem}[The same mistake path for standard softplus]
Replace ReLU by $\phi(t)=\log(1+e^t)$ and use
\[
 (U_j)_i\sim N(0,N\sigma^2/d),\qquad
 a_j\sim N(0,\sigma^2/N)
\]
independently. For each fixed $k\ge2$, there is a finite threshold
$N_\phi(k)$ such that, at every fixed $N\ge N_\phi(k)$, the correct-label
posterior probabilities at the same $T_k$ designated rounds tend to
zero as $\sigma\downarrow0$, with only the observed-margin likelihood (1).
The target has a one-hidden-layer standard-softplus teacher of width
at most $C_k$, margin greater than one, and coefficients bounded by
$8C_k^3$. For an absolute $A>0$ one can take
$N_\phi(k)\le\lceil e^{A k\log(k+1)}\rceil$.
This bound is not numerically optimized, and the explicit ReLU threshold
$N_0$ is not asserted to suffice.
\end{theorem}
\begin{proof}
We first state precisely the additional inherited analytic input.
For any fixed subset $\mathcal T$ of full rail transversals, write
$Q_{\mathcal T}$ for its weighted absolute Fourier sum and
$\mathcal D_{\mathcal T}$ for its downward span.
The bias-independent integral estimate and stability argument in the
softplus section of \cite{shared} give a finite constant $K_{\mathcal T}$
such that
\[
 Q_{\rm full}(\phi(b_0+w\cdot x))\le\|w\|_2,\qquad
 \|\phi(b_0+w\cdot x)-P_{\mathcal D_{\mathcal T}}
                     \phi(b_0+w\cdot x)\|_2^2
 \le K_{\mathcal T}\|w\|_2
            [\|w\|_2-Q_{\mathcal T}(\phi(b_0+w\cdot x))]. \tag{27}
\]
These statements hold for arbitrary bias $b_0$.
The proof first bounds the weighted absolute coefficient integral by
$\|w\|_2$, uniformly in bias. Averaging ReLU against a centered logistic
bias gives softplus. The full rail stability estimate restricts to any
subset by charging a closest omitted ray's coefficient to the omitted
mass; this restriction does not require the subset to be a linear code.
Finally the activation is 1-Lipschitz, and a function of a retained
ray belongs to its downward span, giving (27).

At the current stage let $\mathcal T$ be exactly (21), and keep the
dual $\Lambda$ from (18). Its coefficients equal their positive
weight on $\mathcal T$, and have absolute ratio at most some
$\vartheta<1$ on the other full transversals. Finiteness and the strict
inequalities in Lemma~\ref{lem:coefficients} give this $\vartheta$.
For every softplus atom $h$,
\[
 |\langle\Lambda,h\rangle|
       \le Q_{\mathcal T}(h)+\vartheta
                            [Q_{\rm full}(h)-Q_{\mathcal T}(h)]
       \le\vartheta\|w\|+(1-\vartheta)Q_{\mathcal T}(h).
\]
Put
\[
 K_N(\theta)=\frac dN\sum_j\|U_j\|^2+N\sum_j a_j^2,\quad
 B=\frac{K_N}{2\sqrt d},\quad D=B-\langle\Lambda,F\rangle.
\]
The arithmetic--geometric mean bound and weighted Cauchy--Schwarz
in (27) imply
\[
 D\ge0,\qquad
 \|F-P_{\mathcal D_{\rm ret}}F\|_2^2
                   \le\frac{K_{\mathcal T}}{1-\vartheta}\,BD. \tag{28}
\]
Indeed $B\ge\sum_j|a_j|\|U_j\|\ge\sum_j|a_j|\|w_j\|$, and the
sum of the signed atom deficits is bounded by $D$.
The activation bound $\phi(t)\le|t|+\log2$ also gives
\[
 \|F\|_\infty\le\sqrt{d+1}\,B
                       +(\log2)\sqrt{2\sqrt d\,B}.       \tag{29}
\]

Consider arbitrary finite output tables paired with $0\le B\le C_v+1$,
satisfying the closed observed margins, (28), and the bound (29).
This is a compact set, independent of width.
The dual sign table gives $B\ge\langle\Lambda,F\rangle\ge C_v$.
At equality, (28) puts $F$ in $\mathcal D_{\rm ret}$ and the active
observed margins are exact. The residual and missing-frequency argument
in Proposition~\ref{prop:unique} applies verbatim to such a table;
it used only these properties after (21). Hence equality forces $F=F_v$.
The correct-sign subset consequently has $B\ge C_v+g_v$ for some
$g_v>0$, if nonempty. Energies above the cap are separated already.
This is a correct-sign energy gap uniform in learner width.

To approach the star energy, take a finite minimum-cost ReLU
representation and balance it so that
$d\sum_j\|u_j\|^2+\sum_j c_j^2=2\sqrt d\,C_v$.
At large width use softplus incoming vectors $\sqrt N u_j$,
readouts $c_j/\sqrt N$, and zero readouts on the other neurons.
The energy stays $2\sqrt d\,C_v$ and the outputs tend uniformly to
$F_v$, since $0\le\phi(t)-\relu(t)\le\log2$.
If $\delta_N>0$ bounds the error and tends to zero, multiplying
readouts by $1+4\delta_N$ makes all observed margins strict for large
$N$, because $(1+4\delta_N)(1/2-\delta_N)>1/2$ when
$0<\delta_N<1/4$. The energy still tends to the star energy.
Thus for every sufficiently large fixed width a strictly feasible
open ball lies below the correct-sign energy gap.
The same Gaussian integral comparison proves the posterior conclusion.
Taking a maximum over finitely many stages gives $N_\phi(k)$.
For its quantitative order, (27)--(29) and the proof of (Q5)--(Q6)
give the same uniform lower gap $e^{-A_0k\log(k+1)}$, with the
additional square-root term from (29) absorbed into $H_{\rm cap}$.
The star uses at most $N_0=e^{O(k)}$ nonzero readouts and has balanced
energy polynomial in $k$. Cauchy--Schwarz therefore bounds its
approximation error by $e^{O(k)}/\sqrt N$.
At $N\ge e^{A k\log(k+1)}$ this error and the energy increase from
margin correction are below the uniform correct-sign gap.
Increasing $A$ covers every $k\ge2$.

For the teacher, take the ReLU construction (25), with at most $C_k$
neurons and each coefficient bounded by $C_k$. Replace each ReLU by
$T^{-1}\phi(T\,\cdot)$ with $T=4C_k^2$.
The total uniform error is at most
$C_k^2\log2/T<1/4$. Multiply readouts by two.
The resulting standard-softplus score has margin greater than one,
the same width, and coefficients bounded by $8C_k^3$.
\end{proof}

\section*{Verification and dependency scope}
The finite checkers enclose actual rail costs with integer square roots.
The first checks all 65,536 logical Walsh coefficients at $k=2$.
A second uses exact prefix-permutation symmetry and integer Krawtchouk
transforms to cover every logical coefficient through $k=7$; it checks
all 126 designated stages for $k=2,\ldots,7$, including the 64-stage
instance. They check the control dual, its support and signs,
complementarity, strict mixed coefficient bounds, and exceptional-fiber
combinatorics in exact rational arithmetic for enclosed dyadic costs,
as well as the enclosing inequalities for the actual radical costs.
They do not enumerate physical neural activation chambers.
Fifteen selected Lean lemmas check control algebra, support rigidity,
the missing-character argument, exceptional-set combinatorics, and
teacher margins. Three further lemmas prove the finite exact-penalty
argument. Two exposed-face sparsification lemmas preserve an observed
output and coefficient cost with at most the observation-space dimension
in atoms; the neural atom and dual identifications remain hypotheses.
The analytic atom theorem and path ordering are
inherited inputs. The full new theorem is not claimed to have a Lean
formalization or independent review.

\begin{thebibliography}{9}
\bibitem{shared} ChatGPT-6 Astra, K. H\"anni, and D. Vaintrob,
\emph{Superpolynomial online mistakes for isotropic Bayesian ReLU and
softplus networks}, November 2026. The earlier theorem imposes a
full-cube absolute-margin restriction before observing labels.
The exact rail norm theorem, path ordering, conditional Fourier
orthogonality, and parity-spline teacher used here are from that paper.
\par\smallskip\noindent
\url{https://github.com/kaarelh/AI-works/tree/main/neural-network-generalization/works-in-progress/online-mistakes/isotropic-relu-allwidth}.
\end{thebibliography}
\end{document}
